\subsection{Tác động cột sống}

\begin{longtblr}[
  caption = {Bảng Chỉ Định, Chống Chỉ Định và Biến Chứng trong Tác Động Cột Sống},
  label = {tab:chi_dinh_chong_chi_dinh_tdcs},
]{
  colspec = {X[1.8,l,m] X[4,l,m]},
  hlines,
  vlines,
  row{1} = {font=\bfseries, bg=gray!15, c},
}
Hạng mục & Nội dung chi tiết \\

\textbf{Chỉ định} & 
• Các bệnh lý cơ - cột sống cơ năng.\newline
• Tình trạng đau lưng.\newline
• Cột sống kém/hạn chế vận động.\newline
• Một số trường hợp biến đổi hình thái đốt sống (lồi, lõm, lệch). \\

\textbf{Chống chỉ định tuyệt đối} & 
• Nhiễm trùng cột sống.\newline
• Khối u rõ rệt.\newline
• Gãy xương.\newline
• Rối loạn đông máu nặng. \\

\textbf{Chống chỉ định riêng}\newline\textit{(Thủ thuật Bật)} & 
• Bệnh nhân có dấu hiệu tổn thương rễ thần kinh nặng.\newline
• Vỡ xương.\newline
\textit{*(Yêu cầu người thực hiện phải có kinh nghiệm)*} \\

\textbf{Biến chứng có thể gặp} & 
• Đau tăng thoáng qua.\newline
• Tụ máu tại chỗ.\newline
• Tổn thương rễ thần kinh (hiếm gặp). \\
\end{longtblr}


\begin{longtblr}[
  caption = {Bảng Phân Biệt Hình Thái Đốt Sống (Lồi, Lõm, Lệch)},
  label = {tab:phan_biet_dot_song},
]{
  colspec = {X[1.2,l,m] X[2,l,m] X[2,l,m] X[2,l,m]},
  hlines,
  vlines,
  rowhead = 1,
  row{1} = {font=\bfseries, bg=gray!15, c},
}
Tiêu Chí Phân Biệt & Đốt Sống Lồi & Đốt Sống Lõm & Đốt Sống Lệch \\
Hình Thái Không Gian & Đốt sống nhô lồi ra phía sau so với đường cong sinh lý. & Đốt sống thụt/lõm về phía trước (hướng về phía bụng). & Đốt sống lệch sang một bên (trái hoặc phải) mà không bị lồi hay lõm. \\
Các Biến Thể & Đơn lồi, liên lồi, lồi trên, lồi dưới. & Thường gặp dạng đơn lõm hoặc lõm - lệch. & Đơn lệch, liên lệch, lệch trên, lệch dưới. \\
Đặc Điểm Lớp Cơ & Lớp cơ quanh đốt sống thường \textbf{dày, co cứng}, xơ sợi (hoặc mềm, teo). & Lớp cơ tương ứng thường bị \textbf{teo mỏng}. & Lớp cơ \textbf{bên lệch thường dày, co cứng}; bên đối diện thường bị \textbf{teo mỏng}. \\
Nhiệt Độ Da & Nhiệt độ da vùng bệnh lý thường \textbf{tăng}. & Nhiệt độ da vùng bệnh lý thường \textbf{giảm}. & Bên lệch có nhiệt độ tăng/thay đổi; bên đối diện có nhiệt độ \textbf{thấp}. \\
Biểu Hiện Cảm Giác & Cảm giác đau khu trú \textbf{tăng}. & Cảm giác đau/đối ứng thường \textbf{giảm}. & Bên lệch đau; bên đối diện cảm giác \textbf{giảm}. \\
\end{longtblr}


\begin{longtblr}[
  caption = {Bảng Tổng Hợp 8 Thủ Thuật Tác Động Cột Sống},
  label = {tab:8_thu_thuat_tdcs},
]{
  colspec = {c X[1.2,l,m] X[2.5,l,m] X[2.5,l,m] X[2.5,l,m]},
  hlines,
  vlines,
  row{1} = {font=\bfseries, bg=gray!15, c},
}
STT & Thủ thuật & Kỹ thuật \& Lực thực hiện & Mục tiêu chính / Tác dụng & Lưu ý \& Ứng dụng lâm sàng \\

1 & \textbf{Áp} & Dùng ngón cái hoặc gan bàn tay với diện tiếp xúc nhỏ/lớn tùy thể; lực áp nhẹ $\rightarrow$ giữ vài giây $\rightarrow$ thả, hoặc áp lặp theo nhịp để tạo sóng. & Tạo kích thích diện hẹp hoặc rộng lên đầu gai đốt sống nhằm thay đổi mô xơ và tăng cường tuần hoàn. & Dùng \textbf{áp diện hẹp} cho trường hợp đốt sống lồi. \\

2 & \textbf{Vuốt} & Dùng lực kéo nhẹ theo chiều dọc hoặc chiều chếch. & Làm giãn lớp cơ và tăng cường dẫn lưu tuần hoàn. & Dùng trong \textbf{tiền xử lý} và \textbf{hậu xử lý} trước/sau thủ thuật mạnh để làm dịu mô và đánh giá phản ứng. \\

3 & \textbf{Vê} & Thực hiện động tác xoay tròn, nhỏ. & Xoa bóp trực tiếp lên lớp cơ và làm mềm các mô xơ. & Thường phối hợp với thủ thuật vuốt ở giai đoạn tiền và hậu xử lý. \\

4 & \textbf{Ấn} & Dùng đầu ngón tay hoặc cán khuỷu tay tác động trực tiếp lên một điểm trọng tâm nhỏ. & Làm tan điểm xơ cứng, kích thích thần kinh địa phương để thay đổi cảm giác đau. & Có thể gây đau; cần linh hoạt điều chỉnh lực theo phản ứng của bệnh nhân. \\

5 & \textbf{Đẩy} & Tác dụng lực hướng ra phía sau hoặc hướng ngang. & Di chuyển nhẹ nhàng cấu trúc đốt sống hoặc lớp cơ bị di lệch. & Kết hợp với định hướng để đưa lực chính xác vào lớp cơ mong muốn. \\

6 & \textbf{Kéo} & Tạo lực kéo dọc theo trục của cột sống. & Làm giãn cột sống và giảm áp lực nội đĩa đệm. & Ưu tiên \textbf{kéo nhẹ} đối với đốt sống bị lõm sâu để phục hồi tuần hoàn trước khi áp mạnh. \\

7 & \textbf{Xoay} & Xoay nhẹ hoặc mạnh tùy thuộc vào nhu cầu điều trị. & Phục hồi chuyển động giữa các đốt sống. & Giúp giải tỏa hạn chế vận động khu trú. \\

8 & \textbf{Bật} & Thao tác tăng lực đột ngột (\textit{thrust}). & Giải phóng kẹt khớp hoặc làm thay đổi cấu trúc mô liên kết. & \textbf{Cẩn trọng}: Cần người có kinh nghiệm thực hiện; chống chỉ định khi có dấu hiệu rễ thần kinh nặng hoặc vỡ xương. Dùng \textbf{bật nhẹ} cho đốt sống lệch. \\
\end{longtblr}
x
\begin{longtblr}[
  caption = {Bảng Quy Trình Tiêu Chuẩn 5 Bước Thực Hiện Tác Động Cột Sống},
  label = {tab:quy_trinh_5_buoc_tdcs},
]{
  colspec = {c X[1.5,l,m] X[4,l,m]},
  hlines,
  vlines,
  row{1} = {font=\bfseries, bg=gray!15, c},
}
Bước & Tên bước & Nội dung công việc \\
1 & \textbf{Khám \& Định vị} & Xác định trọng điểm, điểm đối động, loại và thể của đốt sống. \\
2 & \textbf{Chuẩn bị} & Đặt bệnh nhân ở tư thế thoải mái, giải thích thao tác và xin đồng ý. \\
3 & \textbf{Tiền xử lý} & Dùng thủ thuật \textbf{vuốt, vê, áp nhẹ} để làm nóng mô và làm dịu cơ. \\
4 & \textbf{Thao tác chính} & Thực hiện thủ thuật chính theo đúng nguyên tắc \textbf{định điểm, định lực, định hướng}. \\
5 & \textbf{Hậu xử lý} & Vuốt hoặc gõ nhẹ; kiểm tra lại nhiệt độ da, cảm giác đau và khả năng vận động. \\
\end{longtblr}

\subsection{Tác động đốt sống cổ}

\begin{longtblr}[
  caption = {Tổng hợp các loại cơ bao quanh vùng cổ},
  label = {tab:co_vung_co},
]{
    % 1 2 28 7
  colspec = {X[1,c] X[2,l] X[2, l] X[7, l]}, % Định dạng 4 cột: căn giữa, trái, trái, trái
  vlines, % Đường kẻ dọc (bỏ đi nếu bạn thích phong cách booktabs thoáng hơn)
  hlines, % Đường kẻ ngang cho toàn bộ bảng
  rowhead = 1, % Lặp lại dòng tiêu đề khi sang trang mới
  row{1} = {font=\bfseries, bg=gray!10}, % Làm đậm và tô nền xám nhẹ cho tiêu đề
}
% Tiêu đề bảng
STT & Tên Tiếng Việt & Tên Tiếng Anh (Latin) & Nhóm cơ giải phẫu \\

% Thân bảng
\textbf{I} & \SetCell[c=3]{l} \textbf{CÁC CƠ NÔNG VÀ CƠ VÙNG CỔ BÊN (Superficial \& Lateral Muscles)} & & \\
1 & Cơ bám da cổ & Platysma & Cơ nông \\
2 & Cơ ức đòn chũm & Sternocleidomastoid (SCM) & {
    Cơ cổ bên
    \newline
    \includegraphics[width=4cm]{"sternocleidomastoid.png"}
} \\

\textbf{II} & \SetCell[c=3]{l} \textbf{NHÓM CƠ TRÊN MÓNG (Suprahyoid Muscles)} & & \\
3 & Cơ nhị thân (Cơ hai thân) & Digastric & Cơ cổ trước \\
4 & Cơ trâm móng & Stylohyoid & Cơ cổ trước \\
5 & Cơ hàm móng & Mylohyoid & Cơ cổ trước \\
6 & Cơ cằm móng & Geniohyoid & Cơ cổ trước \\

\textbf{III} & \SetCell[c=3]{l} \textbf{NHÓM CƠ DƯỚI MÓNG (Infrahyoid Muscles)} & & \\
7 & Cơ ức móng & Sternohyoid & Cơ cổ trước \\
8 & Cơ ức giáp & Sternothyroid & Cơ cổ trước \\
9 & Cơ giáp móng & Thyrohyoid & Cơ cổ trước \\
10 & Cơ vai móng & Omohyoid & Cơ cổ trước \\

\textbf{IV} & \SetCell[c=3]{l} \textbf{CÁC CƠ BẬC THANG (Scalene Muscles)} & & \\
11 & Cơ bậc thang trước & Anterior scalene & Cơ sâu bên cổ \\
12 & Cơ bậc thang giữa & Middle scalene & Cơ sâu bên cổ \\
13 & Cơ bậc thang sau & Posterior scalene & Cơ sâu bên cổ \\

\textbf{V} & \SetCell[c=3]{l} \textbf{CÁC CƠ TRƯỚC ĐỐT SỐNG (Prevertebral Muscles)} & & \\
14 & Cơ dài cổ & Longus colli & Cơ cổ trước sâu \\
15 & Cơ dài đầu & Longus capitis & Cơ cổ trước sâu \\
16 & Cơ thẳng đầu trước & Rectus capitis anterior & Cơ cổ trước sâu \\
17 & Cơ thẳng đầu bên & Rectus capitis lateralis & Cơ cổ trước sâu \\

\textbf{VI} & \SetCell[c=3]{l} \textbf{NHÓM CƠ GÁY VÙNG CỔ SAU (Posterior Neck Muscles)} & & \\
18 & Cơ thang & Trapezius & {
    Lớp nông
    \newline
    \includegraphics[width=4cm]{"trapezius.png"}
} \\
19 & Cơ gối đầu & Splenius capitis & {
    Lớp giữa
    \newline
    \includegraphics[width=4cm]{"splenius_capitis.png"}
}\\
20 & Cơ gối cổ & Splenius colli (cervicis) & {
    Lớp giữa
    \newline
    \includegraphics[width=4cm]{"splenius_cervicis.png"}
}\\
21 & Cơ nâng vai & Levator scapulae & {
    Lớp giữa
    \newline
    \includegraphics[width=4cm]{"levator_scapulae.png"}
    
} \\
22 & Cơ bán gai đầu & Semispinalis capitis & {
    Lớp sâu
    \newline
    \includegraphics[width=4cm]{"semispinalis_capitis.png"}
} \\
23 & Cơ dài đầu / cổ sau & Longissimus capitis / cervicis & {
    Lớp sâu
    \newline
    \includegraphics[width=4cm]{"longissimus_capitis.png"}
    \includegraphics[width=4cm]{"longissimus_cervicis.png"}
} \\
\textbf{VII} & \SetCell[c=3]{l} \textbf{NHÓM CƠ DƯỚI CHŨM (Suboccipital Muscles)} & & \\
24 & Cơ chéo đầu dưới & Obliquus capitis inferior & {
    Lớp rất sâu
    \newline
    \includegraphics[width=4cm]{"obliquus_capitis_inferior.png"}
} \\
25 & Cơ chéo đầu trên & Obliquus capitis superior & {
    Lớp rất sâu
    \newline
    \includegraphics[width=4cm]{"obliquus_capitis_superior.png"}
} \\
26 & Cơ thẳng đầu sau lớn & Rectus capitis posterior major & {
    Lớp rất sâu
    \newline
    \includegraphics[width=4cm]{"rectus_capitis_posterior_major.png"}
} \\
27 & Cơ thẳng đầu sau nhỏ & Rectus capitis posterior minor & {
    Lớp rất sâu
    \newline
    \includegraphics[width=4cm]{"rectus_capitis_posterior_minor.png"}
} \\
\end{longtblr}


\begin{longtblr}[
  caption = {Tổng hợp các loại xương vùng cổ},
  label = {tab:xuong-vung-co-long},
]{
  colspec = {X[1,c] X[2,l] X[2,l] X[7,l]}, % Định dạng cột: căn lề, canh giữa theo chiều dọc, tỷ lệ độ rộng
  hlines, % Thêm đường kẻ ngang giữa các dòng (nếu muốn bảng thoáng hơn theo chuẩn booktabs thì có thể bỏ dòng này)
  vlines,
  row{1} = {font=\bfseries, bg=gray!10}, % Làm đậm và tô nền nhẹ cho dòng tiêu đề
  rowhead = 1, % Lặp lại dòng tiêu đề nếu bảng bị nhảy sang trang sau
}
STT & Tên tiếng Việt & Tên tiếng Anh & Ghi chú / Đặc điểm chuyên biệt \\
1 & Đốt sống cổ C1 (Đốt đội) & Atlas (First cervical vertebra) & {
        Nâng đỡ hộp sọ, giúp đầu thực hiện chuyển động gật
        \newline
        \includegraphics[width=4cm]{"C1.png"}
    } \\
2 & Đốt sống cổ C2 (Đốt trục) & Axis (Second cervical vertebra) & {
        Có chỏm răng nhô lên làm trục giúp đầu xoay trái/phải
        \newline
        \includegraphics[width=4cm]{"C2.png"}
    } \\
3 & Đốt sống cổ C3 & Cervical vertebra C3 & {
        Đốt sống cổ điển hình
        \newline
        \includegraphics[width=4cm]{"C3.png"}
    } \\
4 & Đốt sống cổ C4 & Cervical vertebra C4 & {
        Đốt sống cổ điển hình 
        \newline
        \includegraphics[width=4cm]{"C4.png"}
    } \\
5 & Đốt sống cổ C5 & Cervical vertebra C5 &{
        Đốt sống cổ điển hình
        \newline
        \includegraphics[width=4cm]{"C5.png"}
    } \\
6 & Đốt sống cổ C6 & Cervical vertebra C6 & {
        Đốt sống cổ điển hình
        \newline
        \includegraphics[width=4cm]{"C6.png"}
    } \\
7 & Đốt sống cổ C7 (Đốt lồi) & Vertebra prominens (C7) & {
        Có gai sống dài nhất, nổi cục rõ phía sau gáy
        \newline
        \includegraphics[width=4cm]{"C7.png"}
    } \\
8 & Xương móng & Hyoid bone & Xương hình chữ U nằm ở trước cổ, không khớp với xương nào khác \\
\end{longtblr}


Đốt sống cổ có bề rộng khoảng 0.5 thốn.
Xác định từ dưới phong phủ, đại chuỳ.
Chia thành 6 khoảng.


\begin{longtblr}[
  caption = {Bảng tổng hợp các kỹ thuật và thủ thuật tác động cột sống vùng cổ - vai},
  entry = {Bảng kỹ thuật Tác động cột sống},
  label = {tab:tdcs}
]{
  colspec = {|X[1,c] |X[2,l] |X[2,l] |X[7,l]|},
  row{1} = {font=\bfseries, bg=gray!10}, % Làm đậm và tô nền nhẹ cho dòng tiêu đề
  hlines,
  vlines,
}
STT & Tên động tác / Thủ thuật & Vùng tác động & Nội dung thực hiện \\
1 & Khám áp lực lớn & Vùng da cột sống cổ - vai & {
    Áp sát bàn tay/ngón tay lên da để nhận biết sự thay đổi nhiệt độ vùng tổn thương.
    \includegraphics[width=\linewidth]{"ap_lon.jpg"}
} \\
2 & Khám áp lực nhỏ & Khối cơ cạnh sống cổ & {
    Dùng lực nhẹ tìm điểm đau khách quan và đánh giá tình trạng co cứng của khối cơ.
    \includegraphics[width=\linewidth]{"ap_nho.jpg"}
} \\
3 & Thủ thuật Vê \newline \textit{(Chẩn bệnh)} & Cạnh sống \newline (các đốt C5, C6, C7...) & {
    Vê nhẹ ngón tay dọc cạnh sống để xác định vị trí đốt sống, phát hiện điểm co cơ/nốt cơ cứng, đồng thời giúp làm mềm cơ tại chỗ.
    \includegraphics[width=\linewidth]{"ve_co.jpg"}
} \\
% 4 & Thủ thuật Lách \newline \textit{(Trị bệnh)} & Khe đốt sống / Huyệt Giáp tích & Lách phần mềm đầu ngón tay vào khe giữa các đốt sống (nơi rễ thần kinh chui ra) để tiếp cận và giải tỏa bờ cơ co cứng. \\
4 & Thủ thuật Xoay \& Tạo sóng & Vùng cột sống cổ & {
    Thực hiện xoay nhẹ nhàng kết hợp phương thức tạo sóng để thư giãn cơ, không dùng lực nén hoặc đẩy mạnh.

} \\
5 & Kỹ thuật Kéo giãn cổ trên & Đốt sống cổ trên \newline (từ C3 trở lên) & Đặt tay ôm/đẩy ở điểm phía trong bả vai (gần đường giữa) rồi thực hiện kéo giãn. \\
6 & Kỹ thuật Kéo giãn cổ dưới & Đốt sống cổ dưới \newline (C4, C5, C6) & Đặt tay ôm/đẩy ở phía bên ngoài vai (gần khớp vai) để thực hiện kéo giãn các đốt phía dưới. \\
\end{longtblr}

% \begin{longtblr}[
%     caption = {Thủ thuật tác động đốt sống cổ},
% ]{
%     width = \textwidth,
%     colspec = {| X[1,c]|X[1, c]|X[1, c]|X[9, l]|},
%     % hlines,
%     rowhead = 1,
% }
%     \hline
%     STT & Kỹ thuật & Số lần & Nội dung
%     \\ \hline
%     1 & Áp &  & {
%         Áp lớn dùng cả lòng bàn tay xác định nhiệt độ \newline
%         Áp nhỏ dùng ngón tay cái và ngón tay trỏ để xác định điểm đau. \newline
%         Vê (dùng ngón tay cái và ngón tay trỏ thực hiện vê) để kiểm tra cạnh sống (khám và làm mềm cơ)\newline
%     }
%     \\ \hline
%     2 & Xoay và tạo sóng & & {
%         Dùng ba đầu ngón tay thực hiện xoay và tạo sóng ở khe đốt sống cổ. \newline
%         % Xoay nhẹ nhàng, lách đầu ngón tay vào đốt sống cổ
%     }
%     \\ \hline
%     3 & Kéo dãn & & {
%         Kéo dãn dùng vai. (Một tay đặt bờ vai, một tay đặt ở điểm cần giãn) (Trên C123) \newline
%         kéo dẫn dùng điểm trong bả vai. (Một tay đặt trong bả vai, một tay đặt ở điểm cần giãn) (dưới C456)
%     }
%     \\ \hline
% \end{longtblr}
